% !TEX root = ../../../main.tex
% 所属：第四部分：Fourier 分析 / 第1章：Fourier级数的基本性质
% 节：1.5 Ces\`aro 和 Abel 求和：Fourier 级数的应用
% 正文按原文件顺序保存；由 ../chapter1.tex 调用。
\section{Ces\`aro 和 Abel 求和：Fourier 级数的应用}

因为 Fourier 级数可能在个别点处发散,为了克服这种缺点,我们寻求通过另一种方式来表示极限
\[
\lim_{N\to\infty}S_N(f)=f.
\]

假定给出一个复数级数
\[
c_0+c_1+c_2+\cdots=\sum_{k=0}^{\infty}c_k.
\]
定义第 $n$ 项和 $s_n$ 为
\[
s_n=\sum_{k=0}^{n}c_k.
\]
\begin{definition}
	我们定义前 $N$ 项和的平均值
\[
\sigma_N=\frac{s_0+s_1+\cdots+s_{N-1}}{N}
\]

量 $\sigma_N$ 被称为级数 $\{s_n\}$ 的 $N$ 次 Ces\`aro 平均,或级数 $\sum\limits_{k=0}^{\infty}c_k$ 的 $N$ 次 Ces\`aro 和。如果当 $N$ 趋于无穷时,$\sigma_N$ 收敛于一个极限 $s$,我们就说级数是 Ces\`aro 可和的,并称 $s$ 为它的 Ces\`aro 和。在函数级数的情形下,我们关心的是级数是逐点收敛的还是一致收敛的。

\end{definition}
我们早先提及 Dirichlet 核不属于好核的类。但它们的平均是性质很好的核。我们用对 Fourier 级数的 $N$ 次 Ces\`aro 和,具体定义如下：
\begin{definition}
	\[
\sigma_N(f)(x)=\frac{S_0(f)(x)+\cdots+S_{N-1}(f)(x)}{N}.
\]
因为 $S_N(f)=f\ast D_N$,故有
\[
\sigma_N(f)(x)=(f\ast F_N)(x),
\]
其中 $F_N(x)$ 是 $N$ 次 Fej\'er 核,定义如下：
\[
F_N(x)=\frac{D_0(x)+\cdots+D_{N-1}(x)}{N}=\sum_{\abs{k}\leqslant N}(1-\frac{|k|}{N})e^{ikx}.
\]
\end{definition}

\begin{lemma}
	可以得到
	\[
	F_N(x)=\frac1N\frac{\sin^2(Nx/2)}{\sin^2(x/2)},
	\]
	且 Fej\'er 核是好核。
\end{lemma}

\begin{proof}
	由定义,
	\[
	\begin{aligned}
		F_N(x)
		&=\frac1N\sum_{k=0}^{N-1}D_k(x)
		=\frac1N\sum_{k=0}^{N-1}\sum_{n=-k}^{k}e^{inx} \\
		&=\frac1N\sum_{|n|\leq N-1}(N-|n|)e^{inx}.
	\end{aligned}
	\]
	另一方面,
	\[
	\begin{aligned}
		\left|\sum_{j=0}^{N-1}e^{ijx}\right|^2
		&=\left(\sum_{j=0}^{N-1}e^{ijx}\right)
		\left(\sum_{k=0}^{N-1}e^{-ikx}\right) \\
		&=\sum_{|n|\leq N-1}(N-|n|)e^{inx}.
	\end{aligned}
	\]
	于是
	\[
	F_N(x)=\frac1N\left|\sum_{j=0}^{N-1}e^{ijx}\right|^2
	=\frac1N\left|\frac{1-e^{iNx}}{1-e^{ix}}\right|^2
	=\frac1N\frac{\sin^2(Nx/2)}{\sin^2(x/2)}.
	\]
	由上式可见 $F_N(x)\geq0$。又因为 $F_N$ 的第 $0$ 个 Fourier 系数为 $1$,所以
	\[
	\frac1{2\pi}\int_{-\pi}^{\pi}F_N(x)\dd x=1.
	\]
	因此
	\[
	\int_{-\pi}^{\pi}|F_N(x)|\dd x=\int_{-\pi}^{\pi}F_N(x)\dd x=2\pi.
	\]
	最后,给定 $\delta>0$,当 $\delta\leq |x|\leq\pi$ 时,$\sin^2(x/2)$ 有正的下界,于是
	\[
	0\leq F_N(x)\leq \frac{C_\delta}{N}.
	\]
	故
	\[
	\int_{\delta\leq |x|\leq\pi}F_N(x)\dd x\to0.
	\]
	所以 $\{F_N\}$ 是好核。
\end{proof}

\begin{theorem}
	如果 $f$ 在圆周上可积,那么 $f$ 的 Fourier 级数在 $f$ 的每个连续点是 Ces\`aro 和于 $f$ 的。并且,如果 $f$ 在圆周上连续,那么 $f$ 的 Fourier 级数是一致 Ces\`aro 和于 $f$ 的。
\end{theorem}

\begin{proof}
	由上一引理,$F_N$ 是好核。又因为
	\[
	\sigma_N(f)(x)=(f\ast F_N)(x),
	\]
	所以由好核逼近定理立即得到结论。
\end{proof}

\begin{corollary}
	若 $f$ 在圆周上可积且对所有 $n$,有 $\widehat f(n)=0$,则在 $f$ 的所有连续点都有
	\[
	f=0.
	\]
\end{corollary}

\begin{proof}
	这显然是成立的。因为 $f$ 的所有部分和都是零,所以其 Ces\`aro 平均也是零。由上一定理,在 $f$ 的连续点处,Ces\`aro 平均收敛到 $f$,于是 $f=0$。
\end{proof}

\begin{corollary}
	圆周上的连续函数可以被三角多项式一致逼近。更精确地说,给定任意 $\eps>0$,存在三角多项式 $P$,使得对任意的 $-\pi\leq x\leq\pi$,
	\[
	|f(x)-P(x)|<\eps
	\]
	成立。
\end{corollary}

\begin{proof}
	由于 Fourier 部分和是三角多项式,所以 Ces\`aro 平均 $\sigma_N(f)$ 也是三角多项式,从而直接由上一定理的一致收敛性得到结论。这个结论就是周期函数的 Weierstrass 逼近定理。
\end{proof}

\begin{definition}
	如果对任意的 $0<r<1$,级数
\[
A(r)=\sum_{k=0}^{\infty}c_k r^k
\]
收敛,且
\[
\lim_{r\to1}A(r)=s,
\]
则复数 $s$ 称为级数 $\sum\limits_{k=0}^{\infty}c_k$ 的 Abel 求和,$A(r)$ 称为此级数的 Abel 平均。如果这个级数收敛于 $s$,那么它是 Abel 和于 $s$
\end{definition}
 另外,Abel 和方法甚至比 Ces\`aro 方法更强。也就是说,当级数是 Ces\`aro 和时,它也是 Abel 和于相同的数。然而,如果考虑级数
\[
1-2+3-4+5-\cdots=\sum_{k=0}^{\infty}(-1)^k(k+1),
\]
由于
\[
A(r)=\sum_{k=0}^{\infty}(-1)^k(k+1)r^k=\frac1{(1+r)^2},
\]
故可以证明它是 Abel 和于 $1/4$,但这个级数不是 Ces\`aro 和的。
\begin{proposition}
	Abel 求和要比一般求和以及 Ces\`aro 求和更强,具体如下
	
	\begin{enumerate}
		\item[(a)] 若复数级数$\sum\limits_{n=1}^{\infty}c_n$
		收敛于一个有限数 $s$,则它的 Abel 和为 $s$
		\item[(b)] 存在 Abel 可求和但是不收敛的级数。
		\item[(c)] 类似地,证明：若$\sum\limits_{n=1}^{\infty}c_n$的 Ces\`aro 和为 $\sigma$,则它的 Abel 和也为 $\sigma$。
		\item[(d)] 给出一个有 Abel 和,但是 Ces\`aro 和发散的例子。
	\end{enumerate}
	这些结果可以总结为
	\[
	\text{收敛}
	\Longrightarrow
	\text{Ces\`aro 可求和}
	\Longrightarrow
	\text{Abel 可求和},
	\]
	其中的箭头均无法逆向成立。
\end{proposition}

\begin{solution}
	\textbf{(a)}
	
	
	先假设$ \sum\limits_{n=1}^{\infty}c_n=0.$令$s_N=\sum\limits_{n=1}^{N}c_n,$则 $s_N\to0$,从而数列 $\{s_N\}$ 有界。
	
	注意到
	\[
	c_n=s_n-s_{n-1},
	\qquad s_0=0,
	\]
	所以
	\[\sum_{n=1}^{N}c_nr^n
	=\sum_{n=1}^{N}(s_n-s_{n-1})r^n
	=\sum_{n=1}^{N}s_nr^n
	-r\sum_{n=1}^{N-1}s_nr^n
	=(1-r)\sum_{n=1}^{N}s_nr^n+s_Nr^{N+1}
	\]
	
	固定 $0<r<1$。由于 $\{s_N\}$ 有界,
	\[
	s_Nr^{N+1}\to0.
	\]
	令 $N\to\infty$,得到
	\[
	\sum_{n=1}^{\infty}c_nr^n
	=(1-r)\sum_{n=1}^{\infty}s_nr^n.
	\]
	
	下面证明右端在 $r\to1^-$ 时趋于 $0$。给定 $\varepsilon>0$,由 $s_n\to0$,存在正整数 $N_0$,使得当 $n>N_0$ 时,
	\[
	|s_n|<\varepsilon.
	\]
	于是
	\[
	\begin{aligned}
		(1-r)\left|\sum_{n=1}^{\infty}s_nr^n\right|
		&\leq
		(1-r)\sum_{n=1}^{N_0}|s_n|r^n
		+(1-r)\sum_{n=N_0+1}^{\infty}|s_n|r^n\\
		&\leq
		(1-r)\sum_{n=1}^{N_0}|s_n|r^n
		+\varepsilon(1-r)\sum_{n=N_0+1}^{\infty}r^n.
	\end{aligned}
	\]
	当 $r\to1^-$ 时,第一项趋于 $0$,而第二项满足
	\[
	\varepsilon(1-r)\sum_{n=N_0+1}^{\infty}r^n
	=\varepsilon r^{N_0+1}
	\leq\varepsilon.
	\]
	因此
	\[
	\lim_{r\to1^-}(1-r)\sum_{n=1}^{\infty}s_nr^n=0.
	\]
	所以
	\[
	\lim_{r\to1^-}\sum_{n=1}^{\infty}c_nr^n=0.
	\]
	
	综上,若
	\[
	\sum_{n=1}^{\infty}c_n=s,
	\]
	则该级数的 Abel 和也是 $s$。
	再说明为什么只需考虑 $s=0$ 的情形。若
	\[
	\sum_{n=1}^{\infty}c_n=s,
	\]
	定义
	\[
	d_1=c_1-s,\qquad d_n=c_n\quad(n\geq2).
	\]
	则
	\[
	\sum_{n=1}^{\infty}d_n=0,
	\]
	并且
	\[
	\sum_{n=1}^{\infty}c_nr^n
	=\sum_{n=1}^{\infty}d_nr^n+sr.
	\]
	因此,只要证明和为 $0$ 的收敛级数的 Abel 和为 $0$,便有
	\[
	\lim_{r\to1^-}\sum_{n=1}^{\infty}c_nr^n
	=0+s=s.
	\]
	
	\textbf{(b)}取	$c_n=(-1)^n.$相应的级数为$-1+1-1+1-\cdots.$其部分和满足
	\[
	s_{2m}=0,\qquad s_{2m-1}=-1,
	\]
	因此该级数不收敛。
	
	但是,当 $0<r<1$ 时,
	\[	\sum_{n=1}^{\infty}c_nr^n
	=\sum_{n=1}^{\infty}(-1)^nr^n
	=-r+r^2-r^3+\cdots
	=-\frac{r}{1+r}.
	\]
	所以
	\[
	\lim_{r\to1^-}\sum_{n=1}^{\infty}c_nr^n
	=-\frac12.
	\]
	因此该级数虽然不收敛,却是 Abel 可求和的,且 Abel 和为$-\frac12.$
	
	事实上,该级数也是 Ces\`aro 可求和的。因为
	\[
	\sigma_N=\frac{s_1+\cdots+s_N}{N},
	\]
	所以
	\[
	\sigma_{2m}
	=\frac{-m}{2m}
	=-\frac12,
	\]
	而
	\[
	\sigma_{2m-1}
	=\frac{-m}{2m-1}
	\to-\frac12.
	\]
	故其 Ces\`aro 和也为 $-\dfrac12$。
	
	\textbf{(c)}设
	\[
	s_n=c_1+\cdots+c_n,
	\qquad
	\sigma_n=\frac{s_1+\cdots+s_n}{n},
	\]
	并假设$\sigma_n\to\sigma.$
	
	类似于第 (a) 问,只需考虑 $\sigma=0$ 的情形。事实上,若定义
	\[
	d_1=c_1-\sigma,\qquad d_n=c_n\quad(n\geq2),
	\]
	则新级数的部分和为$t_n=s_n-\sigma,$其 Ces\`aro 平均为
	\[
	\frac{t_1+\cdots+t_n}{n}
	=\sigma_n-\sigma\to0.
	\]
	同时
	\[
	\sum_{n=1}^{\infty}c_nr^n
	=\sum_{n=1}^{\infty}d_nr^n+\sigma r.
	\]
	因此只需证明 Ces\`aro 和为 $0$ 时,Abel 和也为 $0$。
	
	下面假设$\sigma_n\to0.$由$n\sigma_n=s_1+\cdots+s_n$可得
	\[
	s_n=n\sigma_n-(n-1)\sigma_{n-1},
	\]
	其中规定 $\sigma_0=0$。
	
	于是
	\[\sum_{n=1}^{\infty}s_nr^n
	=\sum_{n=1}^{\infty}
	\bigl(n\sigma_n-(n-1)\sigma_{n-1}\bigr)r^n
	=\sum_{n=1}^{\infty}n\sigma_nr^n
	-r\sum_{n=1}^{\infty}n\sigma_nr^n
	=(1-r)\sum_{n=1}^{\infty}n\sigma_nr^n.
	\]
	由第 (a) 问中的分部求和公式,
	\[
	\sum_{n=1}^{\infty}c_nr^n
	=(1-r)\sum_{n=1}^{\infty}s_nr^n.
	\]
	因此
	\[
	\sum_{n=1}^{\infty}c_nr^n
	=(1-r)^2\sum_{n=1}^{\infty}n\sigma_nr^n.
	\]
	
	给定 $\varepsilon>0$,由 $\sigma_n\to0$,存在正整数 $N_0$,使得当 $n>N_0$ 时,
	\[
	|\sigma_n|<\varepsilon.
	\]
	于是
	\[(1-r)^2
	\left|\sum_{n=1}^{\infty}n\sigma_nr^n\right|
	\leq
	(1-r)^2\sum_{n=1}^{N_0}n|\sigma_n|r^n
	+
	\varepsilon(1-r)^2
	\sum_{n=N_0+1}^{\infty}nr^n.\]
	当 $r\to1^-$ 时,第一项趋于 $0$。又因为
	\[
	\sum_{n=1}^{\infty}nr^n=\frac{r}{(1-r)^2},
	\]
	所以第二项不超过
	\[
	\varepsilon(1-r)^2
	\sum_{n=1}^{\infty}nr^n
	=\varepsilon r
	\leq\varepsilon.
	\]
	故
	\[
	\lim_{r\to1^-}
	(1-r)^2\sum_{n=1}^{\infty}n\sigma_nr^n=0.
	\]
	
	因此,当 Ces\`aro 和为 $0$ 时,Abel 和也为 $0$。由前面的证明可知,若级数的 Ces\`aro 和为 $\sigma$,则其 Abel 和也是 $\sigma$。
	
	\textbf{(d)}取$c_n=(-1)^{n-1}n,$即考虑级数$	1-2+3-4+5-\cdots.$当 $0<r<1$ 时,由
	\[
	\sum_{n=0}^{\infty}(-r)^n=\frac{1}{1+r}
	\]
	两边对 $r$ 求导,得到
	\[
	\sum_{n=1}^{\infty}(-1)^{n-1}nr^{n-1}
	=\frac{1}{(1+r)^2}.
	\]
	两边乘以 $r$,得到
	\[
	\sum_{n=1}^{\infty}(-1)^{n-1}nr^n
	=\frac{r}{(1+r)^2}.
	\]
	因此
	\[
	\lim_{r\to1^-}
	\sum_{n=1}^{\infty}(-1)^{n-1}nr^n
	=\frac14.
	\]
	所以该级数是 Abel 可求和的,且 Abel 和为$\dfrac14.$
	
	
	下面证明它不是 Ces\`aro 可求和的。首先证明：若级数 $\sum c_n$ 是 Ces\`aro 可求和的,则$\dfrac{c_n}{n}\to0.$
	若 $\sigma_n\to\sigma$,则
	\[
	s_n=n\sigma_n-(n-1)\sigma_{n-1},
	\]
	从而
	\[
	\frac{s_n}{n}
	=\sigma_n-\frac{n-1}{n}\sigma_{n-1}
	\to\sigma-\sigma=0.
	\]
	又因为
	\[
	c_n=s_n-s_{n-1},
	\]
	所以
	\[
	\frac{c_n}{n}
	=\frac{s_n}{n}
	-\frac{n-1}{n}\frac{s_{n-1}}{n-1}
	\to0.
	\]
	
	但是对于$c_n=(-1)^{n-1}n,$有	$\dfrac{c_n}{n}=(-1)^{n-1},$
	
	它不趋于 $0$。因此级数$1-2+3-4+5-\cdots$不是 Ces\`aro 可求和的。
	
	综上,
	\[
	\boxed{
		\text{收敛}
		\Longrightarrow
		\text{Ces\`aro 可求和}
		\Longrightarrow
		\text{Abel 可求和}
	}
	\]
	且两个箭头均不能反向成立。
\end{solution}
下面介绍Tauber的一个定理,对上面的$c_n$
,添加一个条件后,上面的箭头就是可逆的
\begin{proposition}
	若$\sum c_n$的Ces\'aro和为$\sigma$,且$c_n = o(\dfrac{1}{n})$,则$\sum c_n$收敛与$\sigma$,对Abel和,有同样的结论
\end{proposition}
\begin{proof}
	考虑\[s_n- \sigma_n =\frac{(n-1)c_n +\cdots+c_2}{n}\]
	由$c_n = o(\dfrac{1}{n})$,$\forall \eps >0,\exists N_0,\text{使得}\forall n>N_0,c_n<\dfrac{\eps}{n}$
	那么
	\[s_n- \sigma_n =\frac{(n-1)c_n +\cdots+c_2}{n} =\frac{c_2+\cdots+(N_0-1)c_{N_0}}{n}+\frac{N_0 c_{N_0+1}\cdots+(n-1)c_n}{n}\]
	对第一项,存在$N>N_0$使,$\forall n>N$,$ \dfrac{c_2+\cdots+(N_0-1)c_{N_0}}{n}<\dfrac{\eps}{2}$对第二项
	\[\frac{N_0 c_{N_0+1}\cdots+(n-1)c_n}{n}<(1-\frac{N_0}{n})\frac{\eps}{2}\],因此
	\[s_n- \sigma_n< c\eps\]
	从而\[\lim\limits_{n\to \infty}s_n- \sigma_n =0\]
	对Abel和,类似的,我们考虑\[\sum_{n=1}^{N}c_n-\sum_{n=1}^{N}c_n r^n\]并且不妨取$r=1-\dfrac{1}{N}$
	\[\sum_{n=1}^{N}c_n-\sum_{n=1}^{N}c_n r^n<\sum_{n=1}^{N}c_n\frac{n}{N}\]类似于上面的方法,不难得出\[\sum_{n=1}^{N}c_n-\sum_{n=1}^{N}c_n r^n\to0\]
\end{proof}


\begin{definition}
	定义函数
\[
f(\theta)\sim\sum_{n=-\infty}^{\infty}a_ne^{in\theta}
\]
的 Abel 平均为
\[
A_r(f)(\theta)=\sum_{n=-\infty}^{\infty}r^{|n|}a_ne^{in\theta}.
\]
\end{definition}
因为指数 $n$ 取正值和负值,很自然地出现 $r^{|n|}$。并且当 $r>0$ 时,$c_n=a_ne^{in\theta}+a_{-n}e^{-in\theta}$,所以 Fourier 级数的 Abel 平均对应于普通级数的 Abel 平均的定义。

因为 $f$ 是可积的,$|a_n|$ 一致有界,所以对每个 $0\leq r<1$,$A_r(f)$ 绝对收敛且一致收敛。正如 Ces\`aro 平均的情形,关键是这些 Abel 平均可以写成卷积的形式,即
\[
A_r(f)(\theta)=(f\ast P_r)(\theta),
\]
其中 $P_r(\theta)$ 是 Poisson 核,定义如下：

	$$P_r(\theta)=\sum_{n=-\infty}^{\infty}r^{|n|}e^{in\theta}.$$

事实上,有
\[
\begin{aligned}
	A_r(f)(\theta)
	&=\sum_{n=-\infty}^{\infty}r^{|n|}a_ne^{in\theta} \\
	&=\sum_{n=-\infty}^{\infty}r^{|n|}
	\left(\frac1{2\pi}\int_{-\pi}^{\pi}f(\varphi)e^{-in\varphi}\dd\varphi\right)e^{in\theta} \\
	&=\frac1{2\pi}\int_{-\pi}^{\pi}f(\varphi)
	\left(\sum_{n=-\infty}^{\infty}r^{|n|}e^{in(\theta-\varphi)}\right)\dd\varphi,
\end{aligned}
\]
其中,交换积分和求和顺序是因为该级数一致收敛。

\begin{lemma}
	如果 $0\leq r<1$,那么有
	\[
	P_r(\theta)=\frac{1-r^2}{1-2r\cos\theta+r^2}.
	\]
	当 $r\to1$ 时,Poisson 核是好核。
\end{lemma}

\begin{proof}
	等式
	\[
	P_r(\theta)=\frac{1-r^2}{1-2r\cos\theta+r^2}
	\]
	已经在前面关于 Poisson 核的例子中给出。注意到由这个表达式可知 $P_r(\theta)>0$。此外,由于 $P_r$ 的第 $0$ 个 Fourier 系数为 $1$,得到
	\[
	\frac1{2\pi}\int_{-\pi}^{\pi}P_r(\theta)\dd\theta=1.
	\]
	因此
	\[
	\int_{-\pi}^{\pi}|P_r(\theta)|\dd\theta=2\pi.
	\]
	最后,给定 $\delta>0$。当 $1/2\leq r\leq1$ 且 $\delta\leq |\theta|\leq\pi$ 时,存在常数 $c_\delta>0$,使得
	\[
	1-2r\cos\theta+r^2\geq c_\delta.
	\]
	于是
	\[
	0\leq P_r(\theta)\leq \frac{1-r^2}{c_\delta}.
	\]
	当 $r\to1$ 时,上式右边趋于 $0$,故
	\[
	\int_{\delta\leq |\theta|\leq\pi}P_r(\theta)\dd\theta\to0.
	\]
	所以当 $r\to1$ 时,Poisson 核构成好核。
\end{proof}

\begin{theorem}
	在圆周上可积函数的 Fourier 级数在每个连续点 Abel 和于 $f$。此外,如果 $f$ 在圆周上连续,那么 $f$ 的 Fourier 级数一致 Abel 和于 $f$。
\end{theorem}

\begin{proof}
	由上一引理,$P_r$ 在 $r\to1$ 时构成好核。又因为
	\[
	A_r(f)(\theta)=(f\ast P_r)(\theta),
	\]
	所以由好核逼近定理可得结论。
\end{proof}
下面,我们考虑$f$在$\theta$处不连续的情况,我们只考虑跳跃点,即第一类间断点
\begin{proposition}

	\begin{enumerate}
		\item[(a)] 若 $f$ 在 $\theta$ 处是跃点,则
		\[
		\lim_{r\to1^-}A_r(f)(\theta)
		=\frac{f(\theta^+)+f(\theta^-)}{2}.
		\]
		
		\item[(b)] 类似的：若 $f$ 在 $\theta$ 处有跳跃间断,
		则 $f$ 在 $\theta$ 处的 Fourier 级数的 Ces\`aro 和为
		\[
		\frac{f(\theta^+)+f(\theta^-)}{2}.
		\]
	\end{enumerate}
	
\end{proposition}


\begin{proof}
	记
	\[
	\alpha=f(\theta^+),\qquad
	\beta=f(\theta^-),\qquad
	L=\frac{\alpha+\beta}{2}.
	\]
	
	\textbf{(a)}
	\[
	A_r(f)(\theta)
	=(P_r*f)(\theta)
	=\frac{1}{2\pi}\int_{-\pi}^{\pi}
	f(\theta-t)P_r(t)\dd t,
	\]
	其中 Poisson 核为
	\[
	P_r(t)=\frac{1-r^2}{1-2r\cos t+r^2}.
	\]
	
	由于 $P_r$ 是偶函数,且
	\[
	\frac{1}{2\pi}\int_{-\pi}^{\pi}P_r(t)\dd t=1,
	\]
	所以
	\[
	\frac{1}{2\pi}\int_{-\pi}^{0}P_r(t)\dd t
	=
	\frac{1}{2\pi}\int_{0}^{\pi}P_r(t)\dd t
	=\frac12.
	\]
	
	在积分区间 $[-\pi,0]$ 上作变量代换 $t\mapsto-t$,并利用
	$P_r(-t)=P_r(t)$,可得
	\[
	\begin{aligned}
		A_r(f)(\theta)
		&=\frac{1}{2\pi}\int_0^\pi
		f(\theta-t)P_r(t)\dd t
		+\frac{1}{2\pi}\int_0^\pi
		f(\theta+t)P_r(t)\dd t\\
		&=\frac{1}{2\pi}\int_0^\pi
		\bigl(f(\theta-t)+f(\theta+t)\bigr)P_r(t)\dd t.
	\end{aligned}
	\]
	
	另一方面,由
	\[
	\frac{1}{2\pi}\int_0^\pi P_r(t)\dd t=\frac12,
	\]
	有
	\[
	L
	=\frac{\alpha+\beta}{2}
	=\frac{1}{2\pi}\int_0^\pi
	(\alpha+\beta)P_r(t)\dd t.
	\]
	因此
	\[
	\begin{aligned}
		A_r(f)(\theta)-L
		&=\frac{1}{2\pi}\int_0^\pi
		\bigl[
		f(\theta+t)-\alpha
		+f(\theta-t)-\beta
		\bigr]P_r(t)\dd t.
	\end{aligned}
	\]
	
	给定 $\varepsilon>0$。由$\lim\limits_{\substack{t\to0\\t>0}}f(\theta+t)=\alpha,
	\lim\limits_{\substack{t\to0\\t>0}}f(\theta-t)=\beta,$
	存在 $\delta>0$,使得当 $0<t<\delta$ 时,
	\[
	|f(\theta+t)-\alpha|<\varepsilon,
	\qquad
	|f(\theta-t)-\beta|<\varepsilon.
	\]
	
	将积分分为 $(0,\delta)$ 和 $(\delta,\pi)$ 两部分。对于第一部分,
	由 $P_r(t)\geq0$,有
	\[\frac{1}{2\pi}\int_0^\delta
		\left|
		f(\theta+t)-\alpha
		+f(\theta-t)-\beta
		\right|P_r(t)\dd t\leq
		\frac{2\varepsilon}{2\pi}
		\int_0^\delta P_r(t)\dd t\leq
		\frac{2\varepsilon}{2\pi}
		\int_0^\pi P_r(t)\dd t
		=\varepsilon.
	\]
	
	对于第二部分,令
	\[
	C_\delta
	=
	\int_\delta^\pi
	\bigl(
	|f(\theta+t)|+|f(\theta-t)|+|\alpha|+|\beta|
	\bigr)\dd t.
	\]
	由于 $f$ 可积,故 $C_\delta<\infty$。于是
	\[\frac{1}{2\pi}\int_\delta^\pi
		\left|
		f(\theta+t)-\alpha
		+f(\theta-t)-\beta
		\right|P_r(t)\dd t\leq
		\frac{C_\delta}{2\pi}
		\sup_{\delta\leq t\leq\pi}P_r(t).\]
	
	当 $\delta\leq t\leq\pi$ 时,
	\[
	1-2r\cos t+r^2
	=(1-r)^2+2r(1-\cos t)
	\geq 2r(1-\cos\delta).
	\]
	因此,当 $r$ 充分接近 $1$ 时,
	\[
	P_r(t)
	\leq
	\frac{1-r^2}{2r(1-\cos\delta)},
	\qquad \delta\leq t\leq\pi.
	\]
	从而
	\[
	\sup_{\delta\leq t\leq\pi}P_r(t)\to0,
	\qquad r\to1^-.
	\]
	
	由此可得
	\[
	\limsup_{r\to1^-}|A_r(f)(\theta)-L|
	\leq\varepsilon.
	\]
	由于 $\varepsilon>0$ 任意,所以
	\[
	\lim_{r\to1^-}A_r(f)(\theta)
	=L
	=\frac{f(\theta^+)+f(\theta^-)}{2}.
	\]

	
	\textbf{(b)}
	由 Fej\'er 平均的卷积表示,
	\[
	\sigma_N(f)(\theta)
	=(F_N*f)(\theta)
	=\frac{1}{2\pi}\int_{-\pi}^{\pi}
	f(\theta-t)F_N(t)\dd t,
	\]
	其中 Fej\'er 核为
	\[
	F_N(t)
	=\frac{1}{N}
	\left(
	\frac{\sin(Nt/2)}{\sin(t/2)}
	\right)^2.
	\]
	
	Fej\'er 核具有下列性质：
	\[
	F_N(t)\geq0,\qquad F_N(-t)=F_N(t),
	\]
	并且
	\[
	\frac{1}{2\pi}\int_{-\pi}^{\pi}F_N(t)\dd t=1.
	\]
	因此
	\[
	\frac{1}{2\pi}\int_0^\pi F_N(t)\dd t=\frac12.
	\]
	
	与第 (a) 问相同,在负半区间作变量代换,得到
	\[
	\sigma_N(f)(\theta)
	=
	\frac{1}{2\pi}\int_0^\pi
	\bigl(f(\theta-t)+f(\theta+t)\bigr)F_N(t)\dd t.
	\]
	同时
	\[
	L
	=
	\frac{1}{2\pi}\int_0^\pi
	(\alpha+\beta)F_N(t)\dd t.
	\]
	所以
	\[
	\begin{aligned}
		\sigma_N(f)(\theta)-L
		&=\frac{1}{2\pi}\int_0^\pi
		\bigl[
		f(\theta+t)-\alpha
		+f(\theta-t)-\beta
		\bigr]F_N(t)\dd t.
	\end{aligned}
	\]
	
	给定 $\varepsilon>0$,取与第 (a) 问相同的 $\delta>0$,使得
	当 $0<t<\delta$ 时,
	\[
	|f(\theta+t)-\alpha|<\varepsilon,
	\qquad
	|f(\theta-t)-\beta|<\varepsilon.
	\]
	
	于是
	\[\frac{1}{2\pi}\int_0^\delta
		\left|
		f(\theta+t)-\alpha
		+f(\theta-t)-\beta
		\right|F_N(t)\dd t
		\leq
		\frac{2\varepsilon}{2\pi}
		\int_0^\delta F_N(t)\dd t\leq
		\frac{2\varepsilon}{2\pi}
		\int_0^\pi F_N(t)\dd t
		=\varepsilon.
	\]
	
	在 $\delta\leq t\leq\pi$ 上
	\[
	F_N(t)
	\leq
	\frac{1}{N\sin^2(\delta/2)}.
	\]
	因此
	\[
	\sup_{\delta\leq t\leq\pi}F_N(t)
	\leq
	\frac{1}{N\sin^2(\delta/2)}
	\to0,
	\qquad N\to\infty.
	\]
	
	利用第 (a) 问中定义的常数 $C_\delta$,有
	\[\frac{1}{2\pi}\int_\delta^\pi
		\left|
		f(\theta+t)-\alpha
		+f(\theta-t)-\beta
		\right|F_N(t)\dd t\leq
		\frac{C_\delta}{2\pi}
		\sup_{\delta\leq t\leq\pi}F_N(t)
		\to0.
	\]
	
	所以
	\[
	\limsup_{N\to\infty}
	|\sigma_N(f)(\theta)-L|
	\leq\varepsilon.
	\]
	由于 $\varepsilon>0$ 任意,得到
	\[
	\lim_{N\to\infty}\sigma_N(f)(\theta)
	=L
	=\frac{f(\theta^+)+f(\theta^-)}{2}.
	\]

\end{proof}
对上面的Tauber定理，Littlewood 给出了 一个细化版本
\begin{proposition}
	\begin{enumerate}
		\item[(a)] 若级数
		\[
		\sum_{n=1}^{\infty}c_n
		\]
		的 Abel 和为 $s$，且
		\[
		c_n=O\left(\frac{1}{n}\right),
		\]
		则
		\[
		\sum_{n=1}^{\infty}c_n
		\]
		在通常意义下收敛于 $s$。
		
		\item[(b)] 从而，若级数
		\[
		\sum_{n=1}^{\infty}c_n
		\]
		的 Ces\`aro 和为 $s$，且
		\[
		c_n=O\left(\frac{1}{n}\right),
		\]
		则
		\[
		\sum_{n=1}^{\infty}c_n
		\]
		在通常意义下收敛于 $s$。
	\end{enumerate}
\end{proposition}
证明极其复杂，这里略过

这些结果可以应用到 Fourier 级数上。若可积函数 $f$
满足
\[
\widehat{f}(\nu)
=O\left(\frac{1}{|\nu|}\right),
\qquad |\nu|\to\infty,
\]
则有：

\begin{enumerate}
	\item[(i)] 若 $f$ 在 $\theta$ 处连续，则当 $N\to\infty$ 时，
	\[
	S_N(f)(\theta)\to f(\theta).
	\]
	
	\item[(ii)] 若 $f$ 在 $\theta$ 处是跃点，则当 $N\to\infty$ 时，
	\[
	S_N(f)(\theta)
	\to
	\frac{f(\theta^+)+f(\theta^-)}{2}.
	\]
	
	\item[(iii)] 若 $f$ 在 $[-\pi,\pi]$ 上连续，则
	\[
	S_N(f)\rightrightarrows f.
	\]
\end{enumerate}
